% Gerada por scripts/tabelas_arranjo_editorial.py; não editar manualmente.
\begingroup
\scriptsize
\setlength{\tabcolsep}{2pt}
\begin{longtable}{lrrrrrrrrr}
\caption[Arranjo sequencial: decisões, revisão e custo]{Arranjo sequencial OpenRecLink e GZ-CMD v3.1, Claude Haiku 5.5, instruções v4: decisões, revisão e custo. São 247 pares verdadeiros na base e 151 em $[5,9)$. PAR~=~1 é definido pelo corte; as classes PAR~=~1 e PAR~=~2 não comprovam incorporação, que depende das evidências usadas na referência. Predição fora da amostra não equivale a validação externa.}
\label{tab:cap8-arranjo-corte}\\
\toprule
Braço & $V_A$ & $S_A$ (\%) & FP automáticos (abaixo) & $V_N$ & $H$ ($V_H$) & LLM & Chamadas & US\$ & $S_H$ (5\%) \\
\midrule
\endfirsthead
\toprule
Braço & $V_A$ & $S_A$ (\%) & FP automáticos (abaixo) & $V_N$ & $H$ ($V_H$) & LLM & Chamadas & US\$ & $S_H$ (5\%) \\
\midrule
\endhead
\bottomrule
\endlastfoot
\multicolumn{10}{l}{\textbf{Vigilância: base inteira}} \\
(i) Corte & 211 & 85,4 & 0 (0) & 36 & 0 (0) & 0 & 0 & 0,0000 & 85,4 \\
Cascata sobre a base inteira & 194 & 78,5 & 34 (34) & 15 & 111 (38) & 347 & 747 & 0,9198 & 93,2 \\
B1 & 194 & 78,5 & 40 (40) & 15 & 115 (38) & 374 & 821 & 1,0178 & 93,2 \\
(ii) Humano integral & 211 & 85,4 & 0 (0) & 0 & 61.485 (teto; 36) & 0 & 0 & 0,0000 & 99,3 \\
\textit{A posteriori} (sem registro interno) & 221 & 89,5 & 34 (34) & 9 & 90 (17) & 330 & 710 & 0,8750 & 96,0 \\
ASR & 215 & 87,0 & 21 (21) & 13 & 96 (19) & 344 & 755 & 0,9380 & 94,4 \\
ASR-Platt & 217 & 87,9 & 3 (3) & 13 & 77 (17) & 110 & 239 & 0,3013 & 94,4 \\
A3 & 219 & 88,7 & 19 (19) & 14 & 100 (14) & 367 & 806 & 1,0107 & 94,0 \\
ASR-PAR2-faixa & 211 & 85,4 & 21 (21) & 12 & 116 (24) & 265 & 583 & 0,7252 & 94,7 \\
\midrule
\multicolumn{10}{l}{\textbf{Vigilância: faixa intermediária da nota final $[5,9)$}} \\
(i) Corte & 123 & 81,5 & 0 (0) & 28 & 0 (0) & 0 & 0 & 0,0000 & 81,5 \\
Cascata sobre a base inteira & 113 & 74,8 & 34 (34) & 12 & 97 (26) & 347 & 747 & 0,9198 & 91,2 \\
B1 & 113 & 74,8 & 40 (40) & 12 & 101 (26) & 374 & 821 & 1,0178 & 91,2 \\
(ii) Humano integral & 123 & 81,5 & 0 (0) & 0 & 21.497 (teto; 28) & 0 & 0 & 0,0000 & 99,1 \\
\textit{A posteriori} (sem registro interno) & 131 & 86,8 & 34 (34) & 8 & 83 (12) & 330 & 710 & 0,8750 & 94,3 \\
ASR & 127 & 84,1 & 21 (21) & 11 & 89 (13) & 342 & 751 & 0,9327 & 92,3 \\
ASR-Platt & 127 & 84,1 & 3 (3) & 12 & 71 (12) & 107 & 232 & 0,2928 & 91,7 \\
A3 & 129 & 85,4 & 19 (19) & 11 & 96 (11) & 362 & 796 & 0,9989 & 92,4 \\
ASR-PAR2-faixa & 123 & 81,5 & 21 (21) & 10 & 109 (18) & 263 & 579 & 0,7199 & 92,8 \\
\midrule
\multicolumn{10}{l}{\textbf{Confirmação: base inteira}} \\
(i) Corte & 211 & 85,4 & 0 (0) & 36 & 0 (0) & 0 & 0 & 0,0000 & 85,4 \\
Cascata sobre a base inteira & 132 & 53,4 & 0 (0) & 28 & 127 (87) & 44 & 95 & 0,1192 & 86,9 \\
B1 & 131 & 53,0 & 4 (4) & 30 & 131 (86) & 70 & 151 & 0,1848 & 86,1 \\
(ii) Humano integral & 211 & 85,4 & 0 (0) & 0 & 61.485 (teto; 36) & 0 & 0 & 0,0000 & 99,3 \\
\textit{A posteriori} (sem registro interno) & 219 & 88,7 & 0 (0) & 17 & 51 (11) & 25 & 53 & 0,0656 & 92,9 \\
ASR & 219 & 88,7 & 4 (4) & 27 & 17 (1) & 41 & 87 & 0,1052 & 89,0 \\
ASR-Platt & 219 & 88,7 & 3 (3) & 27 & 16 (1) & 25 & 54 & 0,0654 & 89,0 \\
A3 & 219 & 88,7 & 5 (5) & 23 & 22 (5) & 32 & 67 & 0,0783 & 90,6 \\
ASR-PAR2-faixa & 217 & 87,9 & 0 (0) & 26 & 22 (4) & 18 & 37 & 0,0441 & 89,4 \\
\midrule
\multicolumn{10}{l}{\textbf{Confirmação: faixa intermediária da nota final $[5,9)$}} \\
(i) Corte & 123 & 81,5 & 0 (0) & 28 & 0 (0) & 0 & 0 & 0,0000 & 81,5 \\
Cascata sobre a base inteira & 65 & 43,0 & 0 (0) & 20 & 106 (66) & 30 & 64 & 0,0793 & 84,6 \\
B1 & 64 & 42,4 & 4 (4) & 22 & 110 (65) & 56 & 120 & 0,1449 & 83,3 \\
(ii) Humano integral & 123 & 81,5 & 0 (0) & 0 & 21.497 (teto; 28) & 0 & 0 & 0,0000 & 99,1 \\
\textit{A posteriori} (sem registro interno) & 127 & 84,1 & 0 (0) & 13 & 51 (11) & 18 & 38 & 0,0471 & 91,0 \\
ASR & 127 & 84,1 & 4 (4) & 23 & 17 (1) & 32 & 68 & 0,0821 & 84,7 \\
ASR-Platt & 127 & 84,1 & 3 (3) & 23 & 16 (1) & 16 & 35 & 0,0423 & 84,7 \\
A3 & 127 & 84,1 & 5 (5) & 19 & 22 (5) & 26 & 54 & 0,0626 & 87,3 \\
ASR-PAR2-faixa & 125 & 82,8 & 0 (0) & 22 & 22 (4) & 9 & 18 & 0,0209 & 85,3 \\
\end{longtable}
\par\footnotesize ASR: arranjo sequencial do registro interno, que preserva o corte do OpenRecLink e aplica a política com recalibração isotônica abaixo dele. PAR~=~2: pares verdadeiros recuperados na revisão manual de origem abaixo do corte. B1: mantém a base inteira e a probabilidade original, substituindo apenas as proporções de erro do revisor. A3: restringe também o treinamento do classificador à região abaixo do corte. ASR-PAR2-faixa: estima a omissão do revisor somente nos pares verdadeiros PAR~=~2 e agrupa as faixas da nota final em $[5,9)$. ASR-Platt: análise de sensibilidade de ASR com \textit{Platt scaling} abaixo do corte. Braço \textit{a posteriori}: cascata sobre a base inteira aplicada abaixo do corte, sem recalibração. Viés de incorporação: uso, na construção da referência, das mesmas evidências avaliadas pelo método.

\noindent\parbox{\linewidth}{\footnotesize $V_A$: pares verdadeiros aceitos automaticamente, inclusive \texttt{MATCH} do modelo de linguagem (LLM); $V_N$: pares verdadeiros perdidos por \texttt{NONMATCH}; $H$: potenciais pares ao revisor humano, com pares verdadeiros $V_H$ entre parênteses. UNSURE do LLM segue ao humano uma única vez. LLM e humano se sobrepõem nesse encaminhamento, sem somar novamente sua carga humana. $S_A=V_A/n$; $S_H=(V_A+0{,}95V_H)/n$ exibe o cenário de sensibilidade humana 0,95, não resultado humano observado. Os três cenários, 0,95, 0,99 e 1,00, constam da Tabela~\ref{tab:cap8-sensibilidade-humana} e de \path{data/cenarios_sensibilidade_humana.csv}; a especificidade humana é perfeita por hipótese. Custos usam somente os registros de consumo dos potenciais pares selecionados por cada braço e faixa, com tarifas históricas por tokens, não faturas. O braço \textit{a posteriori} preserva o corte e aplica a cascata sobre a base inteira abaixo dele, sem recalibração. Fonte: \texttt{data/arranjo\_corte\_metricas.csv} e \texttt{data/arranjo\_corte\_par2.csv}.}
\begin{longtable}{lrrr}
\caption{Partição dos 36 pares verdadeiros PAR~=~2 abaixo do corte. Cada célula mostra a contagem e seu IC de Wilson de 95\%, em porcentagem. As ressalvas sobre a referência operacional e a dependência entre pares permanecem.}
\label{tab:cap8-arranjo-corte-par2}\\
\toprule
Braço & Aceitos (IC 95\%) & Humano (IC 95\%) & Perdidos (IC 95\%) \\
\midrule
\multicolumn{4}{l}{\textbf{Vigilância}} \\
(i) Corte & 0 [0,0; 9,6] & 0 [0,0; 9,6] & 36 [90,4; 100,0] \\
Cascata sobre a base inteira & 10 [15,8; 44,0] & 17 [32,0; 63,0] & 9 [13,8; 41,1] \\
B1 & 8 [11,7; 38,1] & 19 [37,0; 68,0] & 9 [13,8; 41,1] \\
(ii) Humano integral & 0 [0,0; 9,6] & 36 [90,4; 100,0] & 0 [0,0; 9,6] \\
\textit{A posteriori} (sem registro interno) & 10 [15,8; 44,0] & 17 [32,0; 63,0] & 9 [13,8; 41,1] \\
ASR & 4 [4,4; 25,3] & 19 [37,0; 68,0] & 13 [22,5; 52,4] \\
ASR-Platt & 6 [7,9; 31,9] & 17 [32,0; 63,0] & 13 [22,5; 52,4] \\
A3 & 8 [11,7; 38,1] & 14 [24,8; 55,1] & 14 [24,8; 55,1] \\
ASR-PAR2-faixa & 0 [0,0; 9,6] & 24 [50,3; 79,8] & 12 [20,2; 49,7] \\
\midrule
\multicolumn{4}{l}{\textbf{Confirmação}} \\
(i) Corte & 0 [0,0; 9,6] & 0 [0,0; 9,6] & 36 [90,4; 100,0] \\
Cascata sobre a base inteira & 8 [11,7; 38,1] & 11 [18,0; 46,9] & 17 [32,0; 63,0] \\
B1 & 8 [11,7; 38,1] & 11 [18,0; 46,9] & 17 [32,0; 63,0] \\
(ii) Humano integral & 0 [0,0; 9,6] & 36 [90,4; 100,0] & 0 [0,0; 9,6] \\
\textit{A posteriori} (sem registro interno) & 8 [11,7; 38,1] & 11 [18,0; 46,9] & 17 [32,0; 63,0] \\
ASR & 8 [11,7; 38,1] & 1 [0,5; 14,2] & 27 [58,9; 86,2] \\
ASR-Platt & 8 [11,7; 38,1] & 1 [0,5; 14,2] & 27 [58,9; 86,2] \\
A3 & 8 [11,7; 38,1] & 5 [6,1; 28,7] & 23 [47,6; 77,5] \\
ASR-PAR2-faixa & 6 [7,9; 31,9] & 4 [4,4; 25,3] & 26 [56,0; 84,2] \\
\bottomrule
\end{longtable}
\par\footnotesize ASR: arranjo sequencial do registro interno, que preserva o corte do OpenRecLink e aplica a política com recalibração isotônica abaixo dele. PAR~=~2: pares verdadeiros recuperados na revisão manual de origem abaixo do corte. B1: mantém a base inteira e a probabilidade original, substituindo apenas as proporções de erro do revisor. A3: restringe também o treinamento do classificador à região abaixo do corte. ASR-PAR2-faixa: estima a omissão do revisor somente nos pares verdadeiros PAR~=~2 e agrupa as faixas da nota final em $[5,9)$. ASR-Platt: análise de sensibilidade de ASR com \textit{Platt scaling} abaixo do corte. Braço \textit{a posteriori}: cascata sobre a base inteira aplicada abaixo do corte, sem recalibração. Viés de incorporação: uso, na construção da referência, das mesmas evidências avaliadas pelo método.

\endgroup
